% NeuroLingua-Bridge -- ELAMI 2026 (Springer LNCS, runningheads)
% Built on the official LLNCS template (llncs.cls v2.24).
\documentclass[runningheads]{llncs}
\usepackage[T1]{fontenc}
\usepackage{graphicx}
\usepackage{amsmath,amssymb}
\usepackage{booktabs}
\usepackage{multirow}
\usepackage{array}
\usepackage{subcaption}
\usepackage{float}
\usepackage{placeins}
\usepackage{color}
\usepackage{hyperref}
\hypersetup{colorlinks=true,linkcolor=blue,citecolor=blue,urlcolor=blue}
\usepackage{cleveref}
\renewcommand\UrlFont{\color{blue}\rmfamily}
\urlstyle{rm}

% --- Modest spacing to meet the page limit (body font, margins and layout unchanged) ---
\setlength{\textfloatsep}{4pt plus 1pt minus 1pt}
\setlength{\floatsep}{4pt plus 1pt minus 1pt}
\setlength{\intextsep}{4pt plus 1pt minus 1pt}
\setlength{\abovecaptionskip}{2pt}
\setlength{\belowcaptionskip}{0pt}
\usepackage{caption}
\captionsetup{font=small}
\usepackage{etoolbox}
\makeatletter
\renewcommand\section{\@startsection{section}{1}{\z@}%
  {-7\p@ \@plus -1\p@ \@minus -1\p@}{3\p@ \@plus 1\p@ \@minus 1\p@}%
  {\normalfont\large\bfseries\boldmath}}
\renewcommand\subsection{\@startsection{subsection}{2}{\z@}%
  {-5\p@ \@plus -1\p@ \@minus -1\p@}{2\p@ \@plus 1\p@ \@minus 1\p@}%
  {\normalfont\normalsize\bfseries\boldmath}}
\makeatother
% References one step smaller than the \small default, with zero inter-item spacing, to meet the page limit
\AtBeginEnvironment{thebibliography}{\scriptsize}
\makeatletter
\patchcmd{\thebibliography}{\usecounter{enumiv}}{\usecounter{enumiv}\itemsep\z@\parsep\z@}{\typeout{BIBPATCH-OK}}{\typeout{BIBPATCH-FAIL}}
\makeatother

\newcommand{\method}{\textsc{NeuroLingua-Bridge}}
\newcommand{\Lquant}{\mathcal{L}_{\mathrm{quant}}}
\newcommand{\Lclip}{\mathcal{L}_{\mathrm{CLIP}}}
\newcommand{\Lgrl}{\mathcal{L}_{\mathrm{site}}}
\newcommand{\Lproto}{\mathcal{L}_{\mathrm{proto}}}
\newcommand{\Ltok}{\mathcal{L}_{\mathrm{tok}}}

\begin{document}

\title{Diagnosis of Autism Spectrum Disorder using LLMs and Multimodal Brain Connectivity Analysis}
\titlerunning{NeuroLingua-Bridge: ASD Diagnosis via LLMs and Brain Connectivity}
\author{Khalid El Khamlichi \and Fatima Ez-Zahraa Bazay \and Ahmed Drissi El Maliani}
\authorrunning{K. El Khamlichi et al.}
\institute{Faculty of Sciences, Mohammed V University, Rabat, Morocco\\
\email{khalid\_elkhamlichi@um5.ac.ma}}

\maketitle

\begin{abstract}
We present \method{}, a framework for site-invariant Autism Spectrum Disorder (ASD) 
diagnosis that maps resting-state functional MRI (fMRI) into the token space of a large language model (LLM) while remaining invariant to the acquisition site. 
Language-aligned fMRI representations degrade sharply across the seventeen sites 
of ABIDE-I, where scanner and protocol differences dominate the diagnostic signal. 
Our framework trains a text-aligned tokenizer jointly with three domain-generalization 
components, namely a multi-class $K$-way site-adversarial gradient-reversal layer, 
site-balanced dynamic latent prototypes, and episodic Group Distributionally Robust 
Optimization (Group-DRO) for the worst-performing site. Under a strict leave-site-out protocol on ABIDE-I across 17 sites parcellated with the CC200 atlas, \method{} reaches $80.0\%$ accuracy, $81.1\%$ AUC, and $74.1\%$ worst-site accuracy, a 15.9-point improvement over a site-agnostic variant, outperforming supervised and foundation-model baselines despite a stricter evaluation protocol. The code is available 
at \url{https://github.com/khalidelkhamlichi9/ASD-Diagnosis-via-LLMs-and-Brain-Connectivity-}.
\keywords{Autism spectrum disorder \and Domain generalization \and Resting-state 
fMRI \and Large language models \and Site harmonization.}
\end{abstract}

\section{Introduction}
\label{sec:intro}
Resting-state fMRI provides noninvasive information about intrinsic brain function, 
but multi-center analysis is confounded by acquisition-site domain shift, where scanner, 
protocol, and population differences induce distribution changes that are unrelated to 
the diagnostic label. On ABIDE-I~\cite{dimartino2014abide}, Autism Spectrum Disorder 
(ASD) classifiers reach $65$--$75\%$ accuracy within a single site but drop sharply once 
evaluated across its seventeen sites.

Language-aligned tokenizers embed brain activations in the token space of a pretrained large language model~\cite{wei2026fmrilm} but include no site-specific mechanism. A site-agnostic tokenizer of this kind, meaning our framework with the domain-generalization components removed, drops from ${\sim}76\%$ within-site to $68.4\%$ across sites, indicating that site shift, rather than diagnostic difficulty, is the dominant failure mode.

Existing approaches address this problem only partially. Connectivity-based supervised models such as BNT~\cite{kan2022bnt}, BrainGNN~\cite{li2021braingnn}, and 
BrainNetCNN~\cite{kawahara2017brainnetcnn} reach $65$--$69\%$ accuracy but train and 
test within the same site, failing to generalize across scanners. Foundation 
models~\cite{caro2024brainlm,dong2024brainjepa} attain $70$--$73\%$ through large-scale pretraining yet lack explicit site adaptation. Statistical harmonization via ComBat~\cite{fortin2018combat} removes batch effects on extracted features but cannot act on learned representations. Binary Domain-Adversarial Neural Network (DANN) training treats all sites as a single nuisance domain and therefore cannot distinguish the $K{=}17$ heterogeneous scanners of ABIDE-I. None of these methods directly integrates site invariance into a language-aligned diagnostic framework. Rather than treating site shift as a post-hoc preprocessing step, we build site invariance directly into the model. This work addresses four research questions:

\textbf{RQ1}~Does $K$-way multi-class site-adversarial learning reduce site shift better than both the site-agnostic fMRI-LM objective and a binary GRL?

\textbf{RQ2}~Do site-balanced latent prototypes improve class separability once site information is removed? 

\textbf{RQ3}~Does episodic Group-DRO improve worst-site generalization on the hardest unseen scanner? 

\textbf{RQ4}~Does the LLM embedding space provide gains beyond a matched-dimensionality non-linguistic projection, and are these gains consistent across architectures?

To address these questions, we propose \method{}, whose tokenizer is trained jointly with a multi-class $K$-way site 
adversary, site-balanced dynamic prototypes~\cite{vandenoord2017vqvae}, and episodic Group-DRO~\cite{sagawa2020groupdro}, then aligned to frozen LLM embedding spaces via CLIP-style objectives~\cite{radford2021clip}.

\section{Method}
\label{sec:method}
\subsection{Overview and Preprocessing}
\label{sec:method-overview}
The complete \method{} pipeline is shown in \cref{fig:method}. Each fMRI scan is preprocessed via CPAC, parcellated with CC200 ($N{=}200$), interpolated to $\mathrm{TR}{=}2.0$\,s ($T{=}160$), and site-wise normalized. All normalization statistics are computed on training sites only and never leak test-site information.

\begin{figure}[t]
\centering
\includegraphics[width=\textwidth]{figures/fig_llms_paper.png}
\caption{Overview of the \method{} pipeline: 
(a)~fMRI preprocessing with CPAC and structured 
fMRI--text descriptor construction; 
(b)~text-aligned tokenizer with $17$-way 
site-adversarial gradient-reversal layer (GRL) 
and dynamic latent prototypes; 
(c)~frozen-LLM alignment across five backbone
embedding spaces;
(d)~episodic Group-DRO meta-learning for
cross-site generalization;
(e)~clinical output and chatbot interface.}
\label{fig:method}
\end{figure}

\subsection{fMRI-Text Descriptor Construction}
\label{sec:method-desc}
Lacking natural fMRI--text pairs, each scan is converted into structured text by inserting four feature sets into predefined templates: functional connectivity, functional gradients, ICA dynamics, and graph metrics.

\subsection{Stage 1: Text-Aligned fMRI Tokenizer}
\label{sec:method-stage1}
Each sequence is encoded by a transformer encoder into a latent space quantized by a normalized EMA vector quantizer with $K{=}8192$ codes of dimension $d{=}128$~\cite{vandenoord2017vqvae}. The decoder reconstructs the signal:
\begin{equation}
\Lquant = \lVert X - D_\phi(\tilde z)\rVert_2^2 + \mathcal{L}_{\mathrm{commit}}.
\label{eq:quant}
\end{equation}
Using the descriptor pairs (\cref{sec:method-desc}), a CLIP-like contrastive loss~\cite{radford2021clip} pulls matching fMRI--text pairs together while pushing mismatched pairs apart:
\begin{equation}
\Lclip = -\frac{1}{B}\sum_{i=1}^{B}
\log \frac{\exp\!\big(\sigma\,\mathrm{sim}(z_i, z_i^{+})\big)}
{\sum_{j=1}^{B}\exp\!\big(\sigma\,\mathrm{sim}(z_i, z_j^{+})\big)},
\label{eq:clip}
\end{equation}
with cosine similarity $\mathrm{sim}(\cdot,\cdot)$, temperature $\sigma$, and batch size $B$.

\subsection{Multi-Class Site-Adversarial Learning}
\label{sec:method-grl}
The fMRI--text discriminator in fMRI-LM~\cite{wei2026fmrilm} ignores the site $k\in\{1,\dots,K{=}17\}$ and therefore cannot separate site structure. We add a site classifier $D_\psi:\mathbb{R}^C\!\to\!\Delta^{K-1}$ preceded by a gradient-reversal layer $\mathcal{R}_\lambda$~\cite{ganin2016dann}, which acts as the identity in the forward pass and multiplies the gradient by $-\lambda$ in the backward pass,
\begin{equation}
\Lgrl = -\frac{1}{B}\sum_{i=1}^{B}\sum_{k=1}^{K}
\mathbb{1}[s_i{=}k]\,\log\big[D_\psi(\mathcal{R}_\lambda(\bar z_i))\big]_k .
\label{eq:site}
\end{equation}
This yields the minimax problem $\min_{\theta}\max_{\psi}\ \mathcal{L}_{\mathrm{task}}-\lambda\,\Lgrl$: $D_\psi$ is trained to recognize the site, while the GRL forces the encoder $E_\theta$ to remove any discriminative site information.

A standard binary DANN~\cite{ganin2016dann} collapses all 17 scanners into one nuisance label, removing only the mean domain shift but not the heterogeneous per-scanner structure. Our $K{=}17$-way formulation lets the discriminator exploit inter-scanner differences individually, yielding stronger invariance (see ablation in \cref{sec:exp-main}).

\emph{Site-invariance property.} It can be shown (see Prop.~S1) that at any saddle point the latent representation $\bar z$ satisfies $I(\bar z\,;\,s)=0$; this generalizes the two-site setting of~\cite{ganin2016dann} to $K{=}17$. In practice, adversarial training only approaches this optimum, and we report the site-cluster collapse in \cref{fig:tsne} as empirical evidence.

\subsection{Dynamic Latent Prototypes}
\label{sec:method-proto}
To enforce class geometry in the invariant space, we introduce $C=2$ prototypes $\{\mu_c\}\subset\mathbb{R}^{C}$, 
one per class; each token is attracted to its own prototype and repelled from the others:
\begin{equation}
\Lproto = \frac{1}{B}\sum_{i=1}^{B}
\Big[\,\lVert \bar z_i-\mu_{c_i}\rVert_2^2
-\beta\,\textstyle\min_{c\neq c_i}\lVert \bar z_i-\mu_{c}\rVert_2^2\,\Big].
\label{eq:proto}
\end{equation}
The prototypes are not optimizer parameters but evolve by an Exponential Moving Average (EMA)~\cite{vandenoord2017vqvae} of the site-balanced class means so that no single site dominates a prototype:
\begin{equation}
\mu_c \leftarrow \eta\,\mu_c + (1-\eta)\,
\frac{1}{K}\sum_{k=1}^{K}
\frac{\sum_{i:\,c_i=c,\,s_i=k}\bar z_i}{\lvert\{i:c_i=c,s_i=k\}\rvert}.
\label{eq:ema}
\end{equation}
Averaging first over subjects within a site and then over sites prevents any single site from overwhelming a prototype. In Stage-1, the objective is $\mathcal{L}_{\mathrm{Stage\text{-}1}}=
\Ltok+\gamma_p\Lproto$ with $\gamma_p{=}0.10$, where $\Ltok=w_1\Lquant+w_2\Lclip+w_3\Lgrl$ combines the tokenizer terms with learnable softmax-normalized weights $(w_1,w_2,w_3)$.

\subsection{Stage 2: LLM Alignment}
\label{sec:method-stage2}
The tokenizer outputs are aligned to a fixed LLM embedding space with 
fMRI-to-fMRI and fMRI-to-text cosine objectives~\cite{radford2021clip}. 
DeepSeek-Coder-1.3B~\cite{deepseek2024} with a $2048$-d embedding space is the 
primary backbone, selected because its large embedding space yields the best 
prototype separability among the five candidates, as confirmed by the backbone 
study in \cref{sec:exp-backbone}. For the backbone study (RQ4) we rerun 
the full pipeline, all other components fixed, on GPT-2 Medium~\cite{radford2019gpt2} 
and ClinicalT5-large~\cite{raffel2020t5} aligned through their hidden states, and 
on the frozen text embeddings of the API-only models Grok and Gemini, which 
probe embedding spaces rather than trainable backbones.

\subsection{Stage 3: Episodic Meta-Learning 
for Cross-Site Generalization}
\label{sec:method-dro}
The attention-pooling layer is followed by a main ASD classifier head and a secondary sex classifier head with weight $0.3$. Because empirical risk minimization lets 
large sites dominate, we use Group-DRO~\cite{sagawa2020groupdro}, with $\mathcal{L}_k(\theta)$ the loss at site $k$:
\begin{equation}
\min_{\theta}\;\max_{q\in\Delta^{K-1}}\;
\sum_{k=1}^{K} q_k\,\mathcal{L}_k(\theta)
\;=\;\min_{\theta}\;\max_k\;
\mathcal{L}_k(\theta), 
\label{eq:dro}
\end{equation}
where the equality holds because the maximum over $q\in\Delta^{K-1}$ is attained at the vertex placing all weight on the worst-performing site, and the 
adversarial distribution $q$ is updated by exponentiated gradients $q_k\propto q_k\exp(\nu\,\mathcal{L}_k)$, 
placing more weight on harder sites. We retain the checkpoint with the best worst-site source validation accuracy.

\section{Experiments}
\label{sec:exp}
\noindent\textbf{Evaluation protocol.} All models use leave-site-out evaluation where target sites are held out and unseen, 
and training and checkpoint selection use source-site data only. We report results on ABIDE-I~\cite{dimartino2014abide} (17 sites, CC200 atlas) in terms of inter-site performance following \cref{sec:method-overview}, with mean\,$\pm$\,std across $5$ seeds.

\subsection{RQ1--RQ3: Contribution of Each Domain-Generalization Module}
\label{sec:exp-main}
\cref{tab:main} (first block) reports an incremental ablation on DeepSeek-Coder-1.3B. A binary GRL that collapses all 17 sites into a single nuisance domain gains only $+2.5$ accuracy and $+3.1$ worst-site accuracy over the no-DG baseline (\textbf{RQ1}), whereas the $17$-way GRL gains $+6.1$/$+8.6$, showing that multi-class site discrimination is essential for heterogeneous multi-scanner data. Dynamic prototypes (\textbf{RQ2}) add a further $+2.7$/$+4.1$ by restoring class geometry after site information is removed; and episodic Group-DRO (\textbf{RQ3}) reaches $80.0\%$ mean and $74.1\%$ worst-site accuracy ($+11.6$/$+15.9$ total), the largest gains where deployment is hardest.

The ``No LLM'' row replaces Stage~2 with the identity (the DG pipeline operates in the raw tokenizer space without any language alignment). The $+8.0$-point accuracy gap between the No-LLM and LLM-aligned configurations confirms that the LLM embedding space provides structure beyond what the DG modules alone can achieve.

\begin{table}[t]
\centering
\caption{Ablation of DG modules and backbone comparison on ABIDE-I 
(inter-site, $5$ seeds, mean\,$\pm$\,std). First block: sequential inclusion of GRL variants, prototypes, and Group-DRO on DeepSeek-Coder-1.3B. Second block: same full DG pipeline across five embedding spaces. Best per column in bold.}
\label{tab:main}
\resizebox{\linewidth}{!}{%
\setlength{\tabcolsep}{5pt}
\begin{tabular}{lcccc}
\toprule
Configuration / Backbone & Acc $\pm$ std & AUC $\pm$ std & F1 $\pm$ std & W-Site$^\dagger$ $\pm$ std\\
\midrule
\multicolumn{5}{l}{\emph{DG module ablation (DeepSeek-Coder-1.3B)}}\\
\quad No DG (GRL disabled, $\lambda{=}0$) & $68.4${\scriptsize$\pm1.2$} & $69.0${\scriptsize$\pm1.3$} & $68.3${\scriptsize$\pm1.1$} & $58.2${\scriptsize$\pm2.1$}\\
\quad + Binary GRL ($K{=}2$)       & $70.9${\scriptsize$\pm1.1$} & $71.5${\scriptsize$\pm1.2$} & $70.7${\scriptsize$\pm1.0$} & $61.3${\scriptsize$\pm2.0$}\\
\quad + $17$-way GRL ($K{=}17$)    & $74.5${\scriptsize$\pm1.0$} & $75.8${\scriptsize$\pm1.1$} & $74.4${\scriptsize$\pm1.0$} & $66.8${\scriptsize$\pm1.8$}\\
\quad + GRL + Prototypes           & $77.2${\scriptsize$\pm0.9$} & $78.6${\scriptsize$\pm0.9$} & $77.1${\scriptsize$\pm0.8$} & $70.9${\scriptsize$\pm1.5$}\\
\quad + GRL + Proto + G-DRO        & $\mathbf{80.0}${\scriptsize$\pm0.7$} & $\mathbf{81.1}${\scriptsize$\pm0.8$} & $\mathbf{79.9}${\scriptsize$\pm0.7$} & $\mathbf{74.1}${\scriptsize$\pm1.3$}\\
\quad No LLM (full DG, no Stage~2) & $72.0${\scriptsize$\pm1.0$} & $73.5${\scriptsize$\pm1.1$} & $71.8${\scriptsize$\pm1.0$} & $63.0${\scriptsize$\pm1.7$}\\
\midrule
\multicolumn{5}{l}{\emph{Full DG pipeline across backbones / embedding spaces}}\\
\quad Rand. Proj. ($D{=}2048$, no lang.)   & $68.5${\scriptsize$\pm1.2$} & $69.1${\scriptsize$\pm1.3$} & $68.2${\scriptsize$\pm1.2$} & $58.8${\scriptsize$\pm2.0$}\\
\quad GPT-2 Medium~\cite{radford2019gpt2}  & $76.0${\scriptsize$\pm1.1$} & $76.2${\scriptsize$\pm1.2$} & $75.8${\scriptsize$\pm1.1$} & $64.3${\scriptsize$\pm1.9$}\\
\quad ClinicalT5-large~\cite{raffel2020t5} & $77.0${\scriptsize$\pm1.0$} & $77.9${\scriptsize$\pm1.1$} & $77.0${\scriptsize$\pm1.0$} & $66.1${\scriptsize$\pm1.7$}\\
\quad DeepSeek-1.3B~\cite{deepseek2024}    & $\mathbf{80.0}${\scriptsize$\pm0.7$} & $\mathbf{81.1}${\scriptsize$\pm0.8$} & $79.9${\scriptsize$\pm0.7$} & $\mathbf{74.1}${\scriptsize$\pm1.3$}\\
\quad Grok (frozen emb.)$^{\ddagger}$      & $79.0${\scriptsize$\pm0.9$} & $79.6${\scriptsize$\pm0.9$} & $79.0${\scriptsize$\pm0.8$} & $69.8${\scriptsize$\pm1.6$}\\
\quad Gemini (frozen emb.)$^{\ddagger}$    & $80.0${\scriptsize$\pm0.8$} & $80.9${\scriptsize$\pm0.8$} & $\mathbf{80.0}${\scriptsize$\pm0.8$} & $72.5${\scriptsize$\pm1.4$}\\
\bottomrule
\end{tabular}}%
\\[3pt]
{\footnotesize $^\dagger$~The worst site is CMU ($n\!\approx\!10$) in all configurations; the small $n$ explains the wider variance on W-Site. $^{\ddagger}$~Grok-1 and Gemini 1.5 Flash are API-only; we align to their frozen text embeddings, not the models themselves.}
\end{table}

\begin{table}[t]
\centering
\caption{Fair comparison under the same inter-site protocol: each architecture is evaluated without and with LLM alignment (DeepSeek-Coder-1.3B, Stage~2). All models use our leave-site-out splits, GRL, prototypes, and Group-DRO. $\Delta$ is the accuracy gain from adding the LLM. Best overall in bold.}
\label{tab:sota}
\resizebox{\linewidth}{!}{%
\setlength{\tabcolsep}{3.5pt}
\begin{tabular}{l ccc ccc c}
\toprule
 & \multicolumn{3}{c}{Without LLM (DG only)} & \multicolumn{3}{c}{With LLM (+DeepSeek+DG)} & \\
\cmidrule(lr){2-4}\cmidrule(lr){5-7}
Backbone & Acc & AUC & W-Site & Acc & AUC & W-Site & $\Delta$Acc\\
\midrule
ComBat+MLP~\cite{fortin2018combat}   & $59.3$ & $61.6$ & $43.3$ & $74.5$ & $75.8$ & $65.4$ & $+15.2$\\
BrainGNN~\cite{li2021braingnn}       & $71.5$ & $72.8$ & $62.0$ & $75.2$ & $76.5$ & $66.8$ & $+3.7$\\
SWiFT~\cite{kim2023swift}            & $74.8$ & $76.0$ & $66.2$ & $76.8$ & $77.9$ & $68.5$ & $+2.0$\\
\method{} (Ours)                      & $72.0$ & $73.5$ & $63.0$ & $\mathbf{80.0}$ & $\mathbf{81.1}$ & $\mathbf{74.1}$ & $+8.0$\\
\bottomrule
\end{tabular}}
\end{table}

\begin{figure}[t]
\centering
\includegraphics[width=0.72\textwidth]{figures/Analyse t-SNE sans et avec prototypes_1.png}
\caption{t-SNE of latent tokens (DeepSeek-Coder-1.3B). 
Without the GRL (a), tokens cluster by acquisition site. With GRL and prototypes (b), site clusters dissolve and diagnostic groups separate around the learned prototypes $\mu_1,\mu_2$ (stars), confirming $I(\bar z;s)\!\to\!0$. For visual clarity, only the 6 training sites with the most balanced ASD/Control ratio are shown to prevent severe overplotting and preserve the visibility of the class geometry.}
\label{fig:tsne}
\end{figure}


\subsection{RQ4: LLM Contribution and Backbone Generality}
\label{sec:exp-backbone}
To isolate the role of the LLM, \cref{tab:sota} compares each architecture in two configurations: \emph{without LLM} (model + DG components only, no Stage-2 alignment) and \emph{with LLM} (+~DeepSeek-Coder-1.3B alignment), under the same inter-site protocol. All baselines were retrained with our leave-site-out splits, GRL, prototypes, and Group-DRO, ensuring a controlled comparison.

The LLM-aligned variant consistently outperforms the non-aligned one across all architectures: ComBat+MLP gains $+15.2$ pp, BrainGNN $+3.7$ pp, SWiFT $+2.0$ pp, and \method{} $+8.0$ pp. These gains demonstrate that the LLM embedding space provides linguistic-geometric structure that helps class separability beyond what the DG modules alone achieve. The largest gain ($+15.2$ pp for ComBat) occurs when the base representation is weakest, suggesting that LLM alignment compensates most effectively for impoverished feature spaces.

Furthermore, the consistency across five LLM backbones in the second block of \cref{tab:main} (all above the no-DG baseline) confirms that the improvement is a property of the DG + language-alignment paradigm rather than an artifact of one backbone (\textbf{RQ4}). This is visible in \cref{fig:tsne}: site clusters in the raw feature space dissolve under the DG modules, while the classes separate around their prototypes.

\noindent\textbf{Qualitative interface.} \method{} can also present predictions in three conversational modes (\cref{fig:chatbot}), a qualitative demonstration grounded in model predictions and network attributions; the displayed coefficients correspond to the learned Stage-1 weights $(w_1,w_2,w_3)$.

\begin{figure}[t]
\centering
\includegraphics[width=0.56\textwidth]{figures/Deepseek-chatbot.png}
\caption{Example predictive interface based on \method{}: DeepSeek-Coder-1.3B answering one patient in three conversational modes, each grounded in the model's prediction and its network attributions (e.g. an amygdala--prefrontal disconnection cue); a qualitative illustration, not a clinical validation.}
\label{fig:chatbot}
\end{figure}

\FloatBarrier

\section{Discussion}
\label{sec:discussion}
The binary GRL reduces site shift partially ($+2.5$ acc), while the $17$-way GRL captures the heterogeneous per-scanner structure ($+6.1$ acc). Prototypes turn the latent geometry from site-driven to diagnosis-driven (\cref{fig:tsne}); Group-DRO exploits this space to boost generalization on the hardest scanner. The ``No LLM'' ablation and \cref{tab:sota} show that the LLM embedding space adds a further $+8.0$ points over the DG modules alone, consistently across five embedding spaces and four architectures. This confirms the gain is a property of the paradigm, achieved under a strict, deployment-like protocol.

\textbf{Limitations \& Future Work.}
While \method{} achieves robust cross-site generalization, frozen LLMs may constrain domain adaptation, and API-only models (Grok, Gemini) use proxy embeddings. Future work will explore parameter-efficient fine-tuning (LoRA/PEFT), multi-modal integration (sMRI/DTI), and validation on larger cohorts (ABIDE~II~\cite{dimartino2017abide2}, UK Biobank).

\section{Conclusion}
\label{sec:conclusion}
We introduced \method{}, coupling a text-aligned fMRI tokenizer with a $17$-way site-adversarial GRL, site-balanced prototypes, and episodic Group-DRO. Under a strict inter-site protocol on ABIDE-I, it reaches $80.0\%$ accuracy and $74.1\%$ worst-site accuracy ($+15.9$ pp over the no-DG baseline), outperforming both non-aligned and LLM-aligned SOTA methods.

\bibliographystyle{splncs04}
\bibliography{ref}

\end{document}
